% Inicio del contenido del manual de seguridad de radiación
\section*{Transporte de Material Radiactivo}

Únicamente personal autorizado (POE) de nuestra empresa podrá transportar material radiactivo en las unidades debidamente equipadas y destinadas para ello, cuando se transporte material radiactivo por necesidades del servicio y operatividad, se debe seguir la ruta más corta y segura y retornarse mínimo una vez al día a las oficinas generales indicando el lugar donde se encuentran y la ruta que está siguiendo. Si por alguna falla mecánica se tiene que detener la unidad nunca y por ningún motivo se dejará la unidad sola en la vía pública con el material radiactivo, se quedará el auxiliar o el técnico al cuidado de la unidad y una persona deberá ir a conseguir o pedir ayuda mecánica, la unidad deberá ser revisada mecánicamente antes de efectuar algún viaje de más de 100 km de distancia, además de contar con llanta de refacción, gato hidráulico, cables pasa corriente, elementos preventivos, cordón con un longuitud de dos veces el tiro directo de la camioneta para coordinarla en caso necesario y con letreros preventivos de Peligro Radiactivo, extinguir de fuego, caja de herramientas básicas para alguna falla mecánica.

% Consideraciones en caso de averías
\subsection*{En caso de averías}
En caso de que la unidad donde se transporte el material radiactivo sufra alguna falla y no pueda ser reparada en el lugar, por ningún motivo se transportará la fuente en compartimientos ocupados por pasajeros (autobús, ferrocarril, avión, correo). Se deberá llamar a oficinas generales para el envío de otra unidad equipada para el transporte.

% Normas de intensidad de radiación
\subsection*{Intensidad máxima de radiación permitida}
La intensidad máxima permitida fuera del vehículo y en la cabina es de 2 mR/hr.

% Precauciones generales
\subsection*{Almacenamiento de material radiactivo}
La unidad nunca se podrá dejar sola en la vía pública cuando las fuentes se encuentren dentro de este.
Cuando el material radiactivo no esté en uso, las fuentes se deben guardar en un depósito destinado a este fin. El lugar de almacenamiento debe estar adecuadamente protegido y debidamente señalado. Deberá asignarse a una persona competente que lo tenga a su cargo.

% Seguridad del lugar de almacenamiento
\subsection*{Seguridad en los lugares de almacenamiento}
Únicamente el personal autorizado puede introducir o retirar material radiactivo del dicho local. Este lugar debe estar protegido contra toda intrusión.
Los lugares de almacenamiento deben elegirse de modo que el riesgo de inundación sea mínimo, así mismo deberán inspeccionar periódicamente para detectar una posible contaminación.